\documentclass[11pt,a4paper,sans]{moderncv}

\moderncvstyle{classic}
\moderncvcolor{purple}
\moderncvicons{marvosym}

\usepackage[utf8]{inputenc}
\usepackage[T1]{fontenc}
\usepackage{lmodern}
\usepackage[scale=0.85]{geometry}
\setlength{\hintscolumnwidth}{3.4cm}
\urlstyle{same}
\nopagenumbers{}

\name{Bohan}{Lu}
\email{bohan.lu@duke.edu}
\homepage{jerrylvx.github.io}
\extrainfo{\httpslink{github.com/jerrylvx}}

\begin{document}

\makecvtitle

\section{Education}
\cventry{2023--present}{PhD in Electrical and Computer Engineering}{Duke University}{Durham, NC}{expected 2027}{Advisor: Kenneth R. Brown}
\cventry{2019--2022}{MSc in Physics}{Freie Universität Berlin}{Germany}{}{Thesis: \href{https://sites.google.com/view/bohanlu96/research}{\emph{Entanglement dynamics of a many-body localized system}}, advisor: Jens Eisert}
\cventry{2015--2019}{BA in Mathematics and Physics}{Washington University in St.\ Louis}{}{}{Cum Laude, Distinction in Mathematics, High Distinction in Physics. Thesis: \emph{Reduced density operators of finite-sized Heisenberg XXZ chains}}

\section{Publications}
\cvitem{2026}{\textbf{Building codes with transversal CCZ using projective geometry and SAT solvers.} \textbf{B.~Lu}, K.~R.~Brown. Preprint, \href{https://arxiv.org/abs/2610.09341}{arXiv:2610.09341}. Code and data: \href{https://github.com/jerrylvx/gate2code-public}{github.com/jerrylvx/gate2code-public}}
\cvitem{2024}{\textbf{Measuring out quasi-local integrals of motion from entanglement.} \textbf{B.~Lu}, C.~Bertoni, S.~J.~Thomson, J.~Eisert. \emph{Communications Physics} \textbf{7}, 17 (2024). \href{https://doi.org/10.1038/s42005-023-01478-5}{doi:10.1038/s42005-023-01478-5}, \href{https://arxiv.org/abs/2301.01787}{arXiv:2301.01787}}

\section{Research Experience}
\cventry{2023--present}{Doctoral Research}{Duke University}{Brown group}{}{
\begin{itemize}
\item Constructed CSS codes with three logical qubits whose logical CCZ gate is implemented by transversal physical $T/T^\dagger$ gates, at thirteen block lengths from $n=48$ to $n=496$, by combining projective geometry with SAT solving.
\item Proved that no CSS code with three logical qubits, $Z$-distance at least 3, and a quasi-transversal CCZ gate exists below $n=39$.
\item Released the code, data, and verification programs for every construction (\href{https://github.com/jerrylvx/gate2code-public}{gate2code-public}).
\item Designed a dynamical error-correcting code from a minimally non-local hypergraph product code (posters at QIP 2025 and QEC 2025).
\item Developed \href{https://github.com/jerrylvx/stim-floquet}{stim-floquet}, an open-source Python package for Floquet measurement schedules, GF(2) homology checks, and local detector discovery for CSS codes, with Stim circuit output.
\end{itemize}}
\cventry{May--Aug 2026}{Architecture Intern}{Xanadu Quantum Technologies}{remote}{}{Research on decoding for quantum LDPC codes. Supervisor: Ben Baragiola.}
\cventry{Oct 2025--Aug 2026}{Consulting Contractor}{Xanadu Quantum Technologies}{remote}{}{Research on fault-tolerant quantum computing architecture. Supervisor: Rafael Alexander.}
\cventry{Jul--Oct 2022}{Research Intern}{HQS Quantum Simulations GmbH}{Karlsruhe, Germany}{}{Developed and benchmarked a spectral resolution method based on Hermite polynomials for resonant-level models, documented in an internal report. Supervisor: Peter Schmitteckert.}
\cventry{2021--2022}{Master's Thesis Research}{Freie Universität Berlin}{Germany}{}{
\begin{itemize}
\item Proposed a protocol to extract the local integrals of motion of a many-body localized system from its entanglement dynamics, published as the Communications Physics paper above.
\item Ran tensor-network simulations on computing clusters and analyzed quantum-information bounds on entanglement growth.
\end{itemize}}
\cventry{2017--2018}{Greg Delos Fellowship}{Washington University in St.\ Louis}{}{}{Studied the space of reduced density operators of finite-sized XXZ Heisenberg ground states. Computed their convex set with the Bethe ansatz and DMRG, and searched for phase-transition signatures with a support vector machine. Advisor: Alexander Seidel. Talk at APS March Meeting 2019.}
\cventry{2016--2017}{Undergraduate Research}{Washington University in St.\ Louis}{}{}{Gauss--Bonnet theorem for two-dimensional closed orbifolds (advisor: Tang Xiang). Maximum-likelihood search for extended halos around TeV blazars (advisors: Manel Errando, Henric Krawczynski).}

\section{Talks and Posters}
\cvitem{2026}{\emph{Small codes with transversal T implementing a logical CCZ gate.} Poster, QEC 2026.}
\cvitem{2025}{\emph{A dynamical error-correcting code from a minimally non-local hypergraph product code.} With K.~Brown. Posters, QIP 2025 and QEC 2025.}
\cvitem{2023}{\emph{Improved decoding of quantum LDPC codes using neural-network enhanced belief propagation.} With A.~Pesah, J.~Roffe, N.~Raveendran. Posters, QIP 2023 and QEC 2023 (Sydney).}
\cvitem{2022}{\emph{Entanglement dynamics in a many-body localized system.} With C.~Bertoni, S.~Thomson, J.~Eisert. Poster, CRC 183 Summer School 2022.}
\cvitem{2019}{\emph{Reduced density operators of finite-sized Heisenberg XXZ chains.} With A.~Seidel. \href{https://ui.adsabs.harvard.edu/abs/2019APS..MARE07003L/abstract}{Talk}, APS March Meeting 2019.}

\section{Skills}
\cvitem{Agentic research}{Run research through LLM coding agents (Claude Code, Codex): research code, simulations, literature work, and paper revision.}
\cvitem{Agent tooling}{Built paper-dep-mapper (private beta), a command-line tool that gives coding agents source-linked review tasks from a LaTeX paper and runs repeatable checks after each edit. Wrote its agent skills for paper review and citation checking.}
\cvitem{Scientific computing}{Python (NumPy, SciPy, Matplotlib), Stim, Git and GitHub, LaTeX, HPC clusters}

\section{Teaching}
\cvitem{Fall 2024}{Teaching Assistant, ECE 520 Quantum Engineering, Duke University}
\cvitem{Winter 2020}{Teaching Assistant, tensor-network seminar, Freie Universität Berlin}
\cvitem{2016--2018}{Teaching Assistant (Physics 216), grader (Math 4111), and physics help-desk mentor, Washington University in St.\ Louis}

\section{Awards}
\cvitem{2026}{Duke Graduate School Conference Travel Award}

\end{document}
